\section{Интегральная формула Коши для функции и для производной
функции.}\label{ux438ux43dux442ux435ux433ux440ux430ux43bux44cux43dux430ux44f-ux444ux43eux440ux43cux443ux43bux430-ux43aux43eux448ux438-ux434ux43bux44f-ux444ux443ux43dux43aux446ux438ux438-ux438-ux434ux43bux44f-ux43fux440ux43eux438ux437ux432ux43eux434ux43dux43eux439-ux444ux443ux43dux43aux446ux438ux438.}

\textbf{Теорема:} Пусть

\begin{enumerate}
\def\labelenumi{\arabic{enumi}.}
\tightlist
\item
  \(G\) - область на \(\mathbb{C}\) с кусочно-гладкой границей;
\item
  \(\partial G\) ориентирована положительно относительно G;
\item
  \(f(z)\) - аналитическая функция на \(G\);
\item
  \(f(z)\) - непрерывная на \(\overline{G}\). Тогда \(\forall z\in G\)
  справедливо равенство \[
  f(z)= \frac{1}{2\pi i}\int_{\partial G} \frac{f(\xi)\,d\xi}{\xi-z}
  \]
\end{enumerate}

\textbf{Теорема:} Пусть

\begin{enumerate}
\def\labelenumi{\arabic{enumi}.}
\tightlist
\item
  \(G\) - область на \(\mathbb{C}\) с кусочно-гладкой границей;
\item
  \(\partial G\) ориентирована положительно относительно G;
\item
  \(f(z)\) - аналитическая функция на \(G\);
\item
  \(f(z)\) - непрерывная на \(\overline{G}\). Тогда \(\forall z\in G\)
  функция \(f(z)\) имеет производные всех порядков в точке \(z\),
  вычисляемые по формуле: \[
  f^{(n)}(z)= \frac{n!}{2\pi i}\int_{\partial G} \frac{f(\xi)}{(\xi-z)^{n+1}}\,d\xi, \ n=1,2,\dots
  \]
\end{enumerate}
